% Discussion: what the VVUQ tests demonstrate, quality and confidence against the
% external standards, the cost and speed half of the thesis, and comparison to
% prior work with practical insight. Figures 04 and 13 are placed here and
% referenced with \autoref. The through-line is the thesis: assurance, not
% generation, is the decision-bearing work, and standards anchoring makes it
% defensible.

The discussion interprets the record rather than recapitulating it. The
through-line is the thesis stated in the introduction: for a single autonomous
humanoid the assurance process is the decision-bearing work, and binding that
process to external standards is what makes its verdicts defensible. The four
subsections take that claim through the test evidence, the credibility argument,
the cost argument, and the comparison to prior work.

\subsection{What the VVUQ tests demonstrate}

The shape of the test suite is itself an argument. Of the 172 tests, the
ten-gate assurance registry carries 64 and the safety module a further twelve, so
the assurance and safety code together hold more than two-fifths of the budget
and outweigh every behavior module they judge. A suite weighted this way is the
thesis written in test counts: the organisation spent most of its verification
effort on deciding whether a behavior may ship, not on producing the behavior.
The tests also pin the safety semantics each standard requires rather than merely
confirming that functions return. 

The catastrophe-gate tests assert that a hard
stop, an e-stop, or a hand-back is produced at the boundary, matching the IEC
80601-2-77 \S201.x hazardous-situation response and the IEC 60601-1 single-fault
safe state \cite{iec8060127, iec60601}. The balance tests assert the
zero-moment-point margin against the support polygon, the ISO 13482 stability
property for a standing service robot \cite{iso13482}. The autonomy tests score
phase-step concordance between the compiled plan and the annotated reference, the
property UL 4600 and IEEE 7009 require of an autonomous system that must default
to a safe state \cite{ul4600, ieee7009}. And the standards-corpus tests confirm
that every gate resolves its designations through the wired manifest,
so the credibility argument is traceable rather than asserted.

\subsection{Quality and confidence against the external standards}

The decision surface is credible because every gate decision is traceable to a
published designation rather than an ad hoc threshold.
\autoref{tab:pergatestandards} lists the standards each gate resolves at runtime,
and the three credibility questions map cleanly onto them. Verification, the
exact 1.0 pass fraction and the 172-test floor beneath it, answers to ASME V\&V
40-2018 \S8 and IEC 62304 software verification \cite{asmevv40, iec62304}.
Validation, the agreement against the independent references, answers to ASME
V\&V 40 \S9 and to the recorded human review and risk control of ISO 14971
\cite{iso14971}. Uncertainty, the coefficient-of-variation bounds and the
escalate path, answers to NASA-STD-7009A, which treats uncertainty as a
first-class credibility factor rather than an afterthought \cite{nasastd7009a}.

Taken together these map onto the United States Food and Drug Administration
framework for assessing the credibility of computational modeling and simulation,
which asks exactly for a risk-proportionate credibility argument of this kind
\cite{fda2023credibility}. \autoref{fig:matrix} shows the binding as a matrix:
all fifteen governing standards are used by at least one gate or by the
cross-cutting credibility set, and the densest anchoring falls on the three
catastrophe gates and on the workhorse robotic-surgery standard IEC 80601-2-77,
which is where the risk and the required credibility are the highest.

\clearpage

\begin{table}[H]
\caption{Resolved per-gate standards (\texttt{papers/VVUQ-02/codegen/config/standards\_map.yaml} and the manifest);
every decision carries these designations.}
\label{tab:pergatestandards}
\centering
\footnotesize
\begin{tabularx}{\textwidth}{L{6.3cm} Y}
\toprule
\textbf{Gate} & \textbf{Resolved governing standards} \\
\midrule
01 bimanual-handeye-servo & IEC 80601-2-77:2019, ISO 9283:1998, ASME V\&V 40-2018 \\
02 dexterous-finger-force & IEC 80601-2-77:2019, ISO/TS 15066:2016 \\
03 whole-body-balance & ISO 13482:2014, IEC 60601-1 \\
04 autonomous-plan-correctness & UL 4600 (2023), IEEE 7009-2024, IEC 62304 \\
05 instrument-grasp-handover & IEC 80601-2-77:2019, ISO 9283:1998 \\
06 vascular-no-fly-hand & IEC 80601-2-77:2019, ISO 14971:2019 \\
07 bimanual-suturing-anastomosis & IEC 80601-2-77:2019, Fistula Risk Score \\
08 perception-scene-understanding & ASME V\&V 40-2018, IEC 62304 \\
09 shared-or-human-collision & ISO/TS 15066:2016, ISO 10218-1:2011, ISO 13482:2014 \\
10 fault-estop-graceful-degradation & IEC 60601-1, ISO 13849-1:2023, IEEE 7009-2024 \\
\bottomrule
\end{tabularx}
\end{table}

\subsection{The cost and speed half of the thesis}

The thesis claims the assurance work was done faster and less expensively than
conventional verification, and the record supports that claim along four
documented lines. Setup was near zero, because the tree is standard library plus
four small packages with guarded optional imports, so everything core ran
immediately rather than after installing a heavy optional matrix. The work was
one integrated loop, reading code, running it, processing the two large
artifacts, writing documentation, re-running lint, committing, and pushing, with
no context switch between terminal, editor, and version-control client. The
assurance breadth exceeded the shipped tests in a single pass, because the agent
ran the full accept, block, and escalate surface and the full safety surface and
processed the 1790-line comparison and the 1000-row stream structurally. 

And the
record is self-verifying, since every command and its verbatim output is captured
for re-derivation. \autoref{fig:cost} renders this as an explicit and honest cost
bridge; the index is illustrative and grounded in direction only, because the
repository records no dollar figures, and the time anchor is concrete in that the
172-test suite runs in about 1.4 seconds while generation is microsecond-scale.
Honesty requires stating what a conventional provisioned server still does
better: live on-prem language-model intent and judging, real Zenodo deposition,
and physics-backed simulation on a graphics processor all need a provisioned
machine with keys and accelerators that the ephemeral container did not have, so
the cost argument is about the assurance and documentation loop, not about
workloads the container could not run.

\subsection{Comparison to prior work and practical insight}

Positioned against contemporary autonomous systems, the contribution is the
emphasis on assurance over generation for a single embodied surgeon.
\autoref{tab:priorwork} contrasts the present work with deterministic
language-model code generation \cite{trooskens2026compiled}, autonomous
clinical-trial statistical programming \cite{yan2026clinagent}, a validated
generative clinical data pipeline \cite{carlisle2026pipeline}, and the
digital-twin VVUQ survey that frames the field \cite{sel2025vvuq}. Each of those
advances the generation or the validation of software; this paper makes the gate
that decides whether generated behavior may ship the unit of automation. The
practical insight is to be honest about how close to real life this is. It is a
planning and assurance reference, not a robot controller; the platform is a
clearly labelled hypothetical 2030 device; and every number is a simulation
result, with the four-entrant comparison being simulation against simulation.


Within those bounds, the framework nonetheless checks real assurance steps off,
determinism from a fixed seed, traceable external standards, pass-or-fail
verification, validation against independent references, a recorded human review,
and a hand-back-to-human default on escalation, so the field does not have to
re-derive them for each new deliverable.\ The regulatory posture is respectful but
opportunistic: this is the kind of fast, standards-anchored assurance the United
States Food and Drug Administration real-time clinical trials direction invites
\cite{fda2026realtime}, and United States leadership in surgical-humanoid VVUQ
tied to external standards would serve patient safety, efficacy, and trial speed
together.

\begin{table}[H]
\caption{Comparison to prior autonomous and VVUQ work; the present study makes
the ship-decision gate the unit of automation.}
\label{tab:priorwork}
\centering
\footnotesize
\begin{tabularx}{\textwidth}{L{4.7cm} Y Y}
\toprule
\textbf{Prior work} & \textbf{Focus} & \textbf{Contrast with this work} \\
\midrule
Deterministic LLM codegen & reproducible code generation & gates generated behavior not code \\
Autonomous statistical program & precise trial program output & adds risk-proportionate decisions \\
Validated generative pipeline & a validation stage on extraction & makes full VVUQ gate deliverable \\
Digital-twin VVUQ survey & vocabulary and perspective & executes vocabulary on embodiment \\
\bottomrule
\end{tabularx}
\end{table}

\newpage



% TEMPLATE
% \begin{figure}[H]
%      \hspace{-0.7cm}
%      \includegraphics[width=1.1\linewidth, trim={0 0 0 0}, clip]{Images/05-clinical-regulatory-standards-wheel.png}
%     \vspace{-0.85cm}
%     \caption{Clinical and regulatory standards wheel for the wired VVUQ-02 input corpus of fourteen consensus standards and two clinical baselines across six domains.}
%     \label{fig:wheel}
% \end{figure}




\vspace*{-1.25cm}

\begin{figure}[H]
   \hspace{-1.3cm}
     \includegraphics[width=1.125\linewidth, trim={0 0 0 0}, clip]{Images/04-gate-standard-binding-matrix.png}
     \vspace{-0.85cm}
    \caption{Binding matrix of the ten VVUQ gates against the fifteen governing external standards, with every standard anchored by at least one gate.}
    \label{fig:matrix}
\end{figure}

\vspace*{-1.25cm}

\begin{figure}[H]
   \hspace{-0.9cm}
     \includegraphics[width=1.1\linewidth, trim={0 235 0 0}, clip]{Images/13-assurance-cost-assessment.png}
     \vspace{-0.85cm}
    \caption{Assurance cost assessment, autonomous cloud VVUQ versus conventional verification, on an illustrative relative index grounded in direction.}
    \label{fig:cost}
\end{figure}
